% ===================================================================
%  TABELO N-RO 3 — Knabeco kaj juneco.
%  Traduction 2026 d'apres texto/io, controlee sur texto/fr et
%  texto/es. Consigne : voir texto/eo/10-tabelo-01.tex.
%
%  Deux scenes, et l'ido subdivise beaucoup plus que le francais :
%  huit des intertitres n'ont pas de vis-a-vis francais. L'esperanto
%  suit l'ido, comme les autres traductions.
% ===================================================================

%%K t03-apar-1 apar
\VUsaut{3.0mm}
\VUtitre{15.0pt}{60}{TABELO N\textsuperscript{o} 3}
\VUsaut{1.6mm}
\VUtitre{11.4pt}{0}{Knabeco kaj juneco.}
\VUsaut{3.4mm}

%%K t03-apar-2 apar
(La 3-a kaj 4-a tabeloj estas tiel kombinitaj ke ili prezentas en kvar
\VUgras{familiaj scenoj} la esencajn nociojn pri la \VUgras{vetero}, la
\VUgras{aĝo}, la \VUgras{familio}, la \VUgras{vestaro} kaj la \VUgras{stato}.)
\VUsaut{4.0mm}

%%K t03-tit-1 sub
\VUcentre{12.0pt}{0}{\textit{Unua sceno.}}
\VUsaut{3.0mm}

%%K t03-tit-2 sub
\VUpk{10.2pt}{40}{Baptoceremonio en vilaĝo.}
\VUsaut{2.4mm}

%%K t03-tit-3 sub
\VUpk{10.2pt}{40}{La printempo.}
\VUsaut{3.0mm}

%%K t03-c1-01-1 p
1. --- Ni estas en la \VUgras{printempo}, en la monatoj \VUgras{marto},
\VUgras{aprilo} aŭ \VUgras{majo}, en la tagoj de la \VUgras{semajno},
\VUgras{lundo}, \VUgras{mardo}, aŭ \VUgras{merkredo}. Estas \VUgras{la deka
horo} matene.

%%K t03-c1-02-1 p
2. --- La \VUgras{vetero} estas bela, la \VUgras{aero} pura. La suno brilas.
Kelkaj \VUgras{nubetoj} \textsuperscript{(2)} supreniras en la blua
\VUgras{ĉielo} \textsuperscript{(1)}.

%%K t03-c1-03-1 p
3. --- La \VUgras{arboj} apenaŭ havas \VUgras{foliojn}. La maldikaj
\VUgras{poploj} \textsuperscript{(3)} kaj la alta \VUgras{platano}
\textsuperscript{(17)} ĝis nun havas nur burĝonojn.

%%K t03-c1-04-1 p
4. --- La \VUgras{hirundoj} \textsuperscript{(4)} revenas kaj flugetas kun
ĝoja pepado ĉirkaŭ la \VUgras{pinto} \textsuperscript{(7)} de la
\VUgras{sonorilejo} \textsuperscript{(5)} kiu kronas la \VUgras{preĝejon}
\textsuperscript{(6)} de la vilaĝo.

%%K t03-tit-4 sub
\VUsaut{3.4mm}
\VUpk{10.2pt}{40}{La preĝejo.}
\VUsaut{3.0mm}

%%K t03-c1-05-1 p
5. --- Tre simpla estas tiu preĝejo kun sia granda \VUgras{portalo}
\textsuperscript{(13)} kaj dika \VUgras{horloĝo} \textsuperscript{(8)}. La
malgranda \VUgras{montrilo} \textsuperscript{(9)} estas sur la cifero X, ĝi
montras la \VUgras{horojn}. La \VUgras{granda} \textsuperscript{(10)} estas sur
XII (tagmezo); ĝi montras la \VUgras{minutojn}.

%%K t03-c1-06-1 p
6. --- En la sonorilejo, la \VUgras{sonoriloj} \textsuperscript{(11)} gaje
sonadas. Ĉe la supro de la tegmento staras malgranda \VUgras{kruco}
\textsuperscript{(12)} ŝtona.

%%K t03-c1-07-1 p
7. --- La preĝejo havas ne nur grandan \VUgras{portalon}
\textsuperscript{(13)}, ĝi havas ankaŭ flankan pordeton. Tra tiu pordeto sinjoro
\VUgras{paroĥestro} \textsuperscript{(15)} vestita per sia \VUgras{sutano},
eliras el la \VUgras{sakristio} \textsuperscript{(14)} kaj iras al la
\VUgras{paroĥejo} \textsuperscript{(19)}.

%%K t03-c1-08-1 p
8. --- Apud li, jen la \VUgras{laktovendistino} \textsuperscript{(24)} kun sia
\VUgras{azeno} \textsuperscript{(23)}. Ŝi revenas el la urbo, en kiun ŝi
kunportis bonan lakton por la infanoj.

%%K t03-tit-5 sub
\VUsaut{4.0mm}
\VUpk{10.2pt}{40}{La fruktoĝardeno.}
\VUsaut{3.4mm}

%%K t03-c1-09-1 p
9. --- Sinjoro Aliborono (la azeno) estas tute kontenta pri tiu matena kuro.
Li delekte enspiras la aeron balzamitan de la parfumo de la \VUgras{floroj}
\textsuperscript{(20)} kiuj kovras la \VUgras{fruktarbojn}
\textsuperscript{(18)} de la najbara \VUgras{fruktoĝardeno}. Tute rozkoloraj
kaj blankaj estas tiuj \VUgras{persikarboj} \textsuperscript{(18)} kaj
\VUgras{abrikotarboj} \textsuperscript{(19)}, kaj ili esperigas bonan rikolton.

%%K t03-c1-10-1 p
10. --- Dume, la malgrandaj \VUgras{abeloj} \textsuperscript{(22)} de la
\VUgras{abelujo} \textsuperscript{(21)} iras tien rabe kolekti sian dolĉan
\VUgras{mielon} zumante.

%%K t03-c1-11-1 p
11. --- Sed kio okazas? Jen ke la knaboj de la vilaĝo alkuras kaj forkurigas
la \VUgras{anserojn} \textsuperscript{(25)} timigitajn. Tio estas la eliro el
la preĝejo de la sekvantaro post la \VUgras{bapto} de la filo de la notario.

%%K t03-c1-12-1 p
12. --- La malgranda Jakobo, en sia \VUgras{lulilo} \textsuperscript{(26)},
vekiĝas pro la ĝoja sonorado de la sonoriloj, kaj lia fratino Magdalena kiu,
ŝi ankaŭ, volas vidi la baptosekvantaron, petas sian patrinon veni por kombi
ŝian hararon kaj vestigi ŝin ĉe la sojlo de la domo.

%%K t03-c1-13-1 p
13. --- La patrino konsentas. Ŝi metas sian \VUgras{kaptukon}
\textsuperscript{(28)}, sian \VUgras{korsaĵon} \textsuperscript{(29)} kaj sian
\VUgras{antaŭtukon} \textsuperscript{(30)}, kaj ŝi venas sidi sur malgranda
seĝo antaŭ la pordo. Magdalena havas nur sian \VUgras{subjupon}
\textsuperscript{(31)} kaj sian \VUgras{korseton} \textsuperscript{(32)}.

%%K t03-c1-14-1 p
14. --- Kiel feliĉa ŝi estas! Ŝi forgesas siajn ŝatatajn florojn, siajn
\VUgras{diantojn} \textsuperscript{(34)}, sur kiujn sidiĝas la
\VUgras{papilioj} \textsuperscript{(35)}, siajn \VUgras{tulipojn}
\textsuperscript{(36)}, siajn \VUgras{konvalojn} \textsuperscript{(37)}, en
\VUgras{potoj} \textsuperscript{(38)}, sur la \VUgras{muro}
\textsuperscript{(33)}, kaj siajn \VUgras{rozojn} \textsuperscript{(39)} kiuj
formas heĝon tiel belan. Ŝi rigardas la riĉan vestaron de la sinjorinoj de la
eskorto.

%%K t03-c1-15-1 p
15. --- La knaboj de la vilaĝo ne tre atentas la vestaron. Ili sin ĵetas
antaŭen kaj interpuŝiĝas por havi \VUgras{bombonojn} \textsuperscript{(55)} aŭ
\VUgras{monerojn} \textsuperscript{(52)}. Unu el ili ploras
\textsuperscript{(40)}.

%%K t03-c1-16-1 p
16. --- La alia paŝis sur la piedon de la \VUgras{hundo}
\textsuperscript{(42)} kiu kriante forkuras kaj forkurigas la
\VUgras{kokinon} \textsuperscript{(43)} kaj ĝiajn \VUgras{kokidojn}
\textsuperscript{(44)}.

%%K t03-tit-6 sub
\VUsaut{5.0mm}
\VUpk{10.2pt}{40}{La bapta defilo.}
\VUsaut{4.0mm}

%%K t03-c1-17-1 p
17. --- Antaŭ la eskorto marŝas la \VUgras{nutristino} \textsuperscript{(45)}.
Ŝi havas belan \VUgras{kufon} \textsuperscript{(46)} kun longaj rubandoj. Ŝi
portas la \VUgras{novnaskiton} \textsuperscript{(47)} tute vestitan per
\VUgras{puntoj} \textsuperscript{(48)}.

%%K t03-c1-18-1 p
18. --- Malantaŭ la nutristino venas la \VUgras{baptopatro}
\textsuperscript{(49)} kaj la \VUgras{baptopatrino} \textsuperscript{(53)}. Ili
disdonas la bombonojn kaj monerojn. La baptopatro havas \VUgras{gantojn}
\textsuperscript{(51)} kaj \VUgras{cilindran ĉapelon} \textsuperscript{(50)}.
La baptopatrino mane tenas belan \VUgras{bukedon} \textsuperscript{(54)}.

%%K t03-c1-19-1 p
19. --- Malantaŭ ili, ni vidas la \VUgras{onklon} \textsuperscript{(56)} kaj la
\VUgras{onklinon} \textsuperscript{(58)} de la infano. La onklino havas
elegantan \VUgras{ĉapelon} \textsuperscript{(59)}, la onklo nigran
\VUgras{redingoton} \textsuperscript{(57)} kaj blankan veŝton. Jen poste la
\VUgras{kuzo} \textsuperscript{(60)} kaj la \VUgras{kuzino}
\textsuperscript{(61)}. Ĉi tiu tenas per la mano la \VUgras{fratinon}
\textsuperscript{(62)} de la novnaskito, la malgrandan Bertan.

%%K t03-c1-20-1 p
20. --- La alia \VUgras{frato} \textsuperscript{(63)} de Berta marŝas
malantaŭe. Li donas la manon al \VUgras{parencino} \textsuperscript{(64)}. La
\VUgras{amikoj} \textsuperscript{(65)} de la familio venas poste.

%%K t03-tit-7 sub
\VUsaut{4.6mm}
\VUpk{10.2pt}{40}{La spektantoj.}
\VUsaut{3.4mm}

%%K t03-c1-21-1 p
21. --- Apud la sekvantaro, jen \VUgras{almozulo} \textsuperscript{(66)}
apogata sur sia \VUgras{lambastono} \textsuperscript{(67)}. Li petas almozon.

%%K t03-c1-22-1 p
22. --- Aliflanke estas malgranda knabo tre \VUgras{ĝentila}
\textsuperscript{(68)}. Li salutas kaj deziras « \VUgras{bonan tagon} » al la
personoj kiujn li konas.

%%K t03-c1-23-1 p
23. --- La vilaĝanoj eliris el siaj hejmoj por vidi la baptan defilon. Ni
vidas \VUgras{kamparanon} \textsuperscript{(69)} kun sia \VUgras{kotona
ĉapeto} kaj apude \VUgras{kamparaninon} \textsuperscript{(70)}. Poste
\VUgras{maljunulinon} \textsuperscript{(71)} apogatan sur sia \VUgras{bastono}
\textsuperscript{(72)}. Fine la \VUgras{muelisto} \textsuperscript{(73)} kun
sia granda \VUgras{felta ĉapelo} \textsuperscript{(74)} kaj juna kamparanino
piedvestita per \VUgras{lignaj ŝuoj} \textsuperscript{(75)}, kiu instruas
marŝi al sia infaneto.

%%K t03-c1-24-1 p
24. --- Tiu sceno estas tre amuza.
\VUsaut{4.0mm}

%%K t03-tit-8 sub
\VUcentre{12.0pt}{0}{\textit{Dua sceno.}}
\VUsaut{3.0mm}

%%K t03-tit-9 sub
\VUpk{10.2pt}{40}{Publika festo.}
\VUsaut{2.4mm}

%%K t03-tit-10 sub
\VUpk{10.2pt}{40}{Navigaj ludoj kaj regatoj.}
\VUsaut{3.0mm}

%%K t03-c2-01-1 p
1. --- La \VUgras{somero} venis. La varmega suno de la monato \VUgras{junio}
(julio aŭ \VUgras{aŭgusto}) instigas la junulojn bani sin kaj boatveturi.

%%K t03-c2-02-1 p
2. --- Sur la dekstra bordo de la fluo, la boatistoj senvestigas sin por
partopreni en la navigaj ludoj kaj regatoj.

%%K t03-c2-03-1 p
3. --- Unu el ili demetas siajn \VUgras{ŝuojn} \textsuperscript{(1)}; li
malligas (malbutonumas) ilin. Liaj du kamaradoj jam estas pretaj. De nun ili
havas nur sian \VUgras{trikotan veston} \textsuperscript{(2)} kaj
\VUgras{banpantaloneton} \textsuperscript{(3)}. Ili konversacias atendante
siajn amikojn. Ili parolas pri la foiro kaj la festo kiu okazas sur la alia
bordo.

%%K t03-c2-04-1 p
4. --- Sur la tero, apud ili, kuŝas la \VUgras{vestoj} de iliaj kamaradoj:
\VUgras{paro} da \VUgras{botetoj} \textsuperscript{(4)}, \VUgras{ŝuoj}
\textsuperscript{(5)}, \VUgras{ŝtrumpetoj} \textsuperscript{(6)},
\VUgras{feltaj} \textsuperscript{(7)} kaj \VUgras{pajlaj} \textsuperscript{(8)}
ĉapeloj kaj \VUgras{kravato} \textsuperscript{(9)}.

%%K t03-c2-05-1 p
5. --- Unu el la junuloj zorge faldis sian \VUgras{jakon}
\textsuperscript{(10)}. Li metas ĝin sur tualettukon, en kiun lia najbaro
baldaŭ enfermos la du \VUgras{vestojn}.

%%K t03-c2-06-1 p
6. --- Sesa knabo malfruas. Li havas ankoraŭ sian \VUgras{kaskedon}
\textsuperscript{(11)} sur la kapo. Li demetis nek sian \VUgras{ĉemizon}
\textsuperscript{(12)}, nek sian \VUgras{kolumon} \textsuperscript{(13)}, nek
siajn \VUgras{manumojn} \textsuperscript{(14)}. Li mane tenas sian
\VUgras{veŝton} \textsuperscript{(17)} kaj havas ankoraŭ sian
\VUgras{pantalonon} \textsuperscript{(16)} kaj \VUgras{ŝelkojn}
\textsuperscript{(15)}. Certe li ne estos preta por la venonta kuro.

%%K t03-c2-07-1 p
7. --- Apud la junuloj, vojeto sur kies flankoj estas multaj \VUgras{kardoj}
\textsuperscript{(18)}, kondukas al la bordo de la rivero, kie la patro Jakobo
preparas inter la \VUgras{kanoj} \textsuperscript{(19)} la
\VUgras{piroteknikaĵon} \textsuperscript{(20)} kiun oni pafos vespere.

%%K t03-tit-11 sub
\VUsaut{4.4mm}
\VUpk{10.2pt}{40}{La kuro.}
\VUsaut{3.4mm}

%%K t03-c2-08-1 p
8. --- Sur la rivero, kelkaj sinjorinoj, ŝirmante sin per siaj ombreloj,
promenas per \VUgras{barko} \textsuperscript{(21)}. Ili sekvas la kuron kiu
estas tre interesa. En ĉiu \VUgras{jolo} \textsuperscript{(22)}, kvar
\VUgras{remistoj} \textsuperscript{(24)} remas per siaj tutaj fortoj, dum la
\VUgras{stiristo} \textsuperscript{(26)} tenas la \VUgras{stirilon}
\textsuperscript{(25)} kaj direktas la fragilan remboaton. La \VUgras{remiloj}
\textsuperscript{(23)} kadence frapas la akvon kaj la malpezaj boatoj navigas
kun kapturniga rapideco.

%%K t03-c2-09-1 p
9. --- Kelkaj junuloj interluktas, apud la pilastro de la ponto, pri la premio
de la \VUgras{ŝmirita masto} \textsuperscript{(28)}. Unu el ili prudente marŝas
sur la fleksebla kaj glitiga masto por atingi la malgrandan \VUgras{flagon}
\textsuperscript{(29)} fiksitan ĉe la ekstremo. Lia malfeliĉa
\VUgras{konkuranto} \textsuperscript{(30)} falis en la akvon. Li naĝas por
alteriĝi.

%%K t03-c2-10-1 p
10. --- Pli aĝaj junuloj \VUgras{(boat)turniras} \textsuperscript{(32)} apud
la \VUgras{kajo} \textsuperscript{(33)}. Al tiuj ĉiuj ludoj ĉeestas
grandnombra \VUgras{publiko} \textsuperscript{(31)} apogata ĉe la
\VUgras{parapeto} de la \VUgras{ponto} \textsuperscript{(27)} kaj amasigita
sur la \VUgras{bordo}. Kelkaj spektantoj, tamen, sin turnas por sekvi la ludon
de la \VUgras{premimasto} \textsuperscript{(45)} aŭ por admiri la
\VUgras{urbestran eskorton} kiu venas por la \VUgras{disdono} de la premioj.

%%K t03-c2-11-1 p
11. --- Unue, jen kelkaj \VUgras{buboj} \textsuperscript{(35)} antaŭirantaj la
\VUgras{muzikistojn} \textsuperscript{(36)}; poste venas la \VUgras{urbestro}
\textsuperscript{(37)}, la adjunktoj kaj la \VUgras{komunuma konsilantaro} de
la urbeto. Fine marŝas la \VUgras{fajrobrigadanoj} \textsuperscript{(38)}.

%%K t03-tit-12 sub
\VUsaut{5.0mm}
\VUpk{10.2pt}{40}{La foiro.}
\VUsaut{3.6mm}

%%K t03-c2-12-1 p
12. --- Grandnombraj homoj estas sur la \VUgras{foirplaco}
\textsuperscript{(39)}. Oni venas vidi la diversajn \VUgras{barakojn}: la
lignajn \VUgras{ĉevalojn} \textsuperscript{(40)}, la \VUgras{cirkon}
\textsuperscript{(41)} ĉe kiu jam komenciĝis la prezentado; la \VUgras{publikan
balon} \textsuperscript{(42)}, kie la junuloj kaj la junulinoj de la urbo ĝuas
la plezuron de la \VUgras{dancado}.

%%K t03-c2-13-1 p
13. --- En la fundo, ni vidas la \VUgras{arenon} \textsuperscript{(44)} ĉe kiu
la brava hispana \VUgras{matadoro} luktas kontraŭ la furioza \VUgras{virbovo}
\textsuperscript{(45)}, poste la \VUgras{relglitejon} \textsuperscript{(49)},
la \VUgras{pafejon} \textsuperscript{(46)} kie oni ekzercas sin uzi la
\VUgras{pafilojn} kaj pafi al la \VUgras{celtabulo}.

%%K t03-c2-14-1 p
14. --- Proksime, jen la \VUgras{foira teatro} \textsuperscript{(47)} en kiu
oni ludas la \VUgras{komedion} kaj la \VUgras{dramon}. Tre proksime, estas la
barako de la \VUgras{giganto} \textsuperscript{(48)} kaj de la \VUgras{nano}.
La \VUgras{staturo} de la unua estas du metroj plus dek kvin centimetroj; la
dua estas apenaŭ tiel granda kiel infano.

%%K t03-c2-15-1 p
15. --- Fine, jen la granda \VUgras{menaĝerio} \textsuperscript{(50)} kun siaj
\VUgras{sovaĝaj bestoj}, sia \VUgras{tigro} \textsuperscript{(51)} kun potencaj
\VUgras{ungegoj}, sia \VUgras{leono} \textsuperscript{(52)} kun densa
\VUgras{kolkrinaro}, siaj du \VUgras{ĝirafoj} \textsuperscript{(53)} kaj sia
\VUgras{elefanto} \textsuperscript{(54)} kiu tiel lerte uzas sian
\VUgras{rostron}.

%%K t03-c2-16-1 p
16. --- La \VUgras{dresisto} \textsuperscript{(55)} kiu dresas tiujn bestojn
estas ĉe la pordo de la barako.
\VUsaut{6.0mm}
\VUfilet{22mm}
